\documentclass[12pt]{article}
\usepackage[margin=0.58in,top=0.62in,bottom=0.48in]{geometry}
\usepackage{lmodern}
\usepackage{microtype}
\usepackage{fancyhdr}
\usepackage{titlesec}
\usepackage{enumitem}
\usepackage{lastpage}
\usepackage[hidelinks]{hyperref}

\setlength{\parindent}{0pt}
\setlength{\parskip}{0.18em}
\setlength{\headheight}{26pt}
\raggedbottom
\titleformat{\section}{\normalsize\bfseries\scshape}{}{0pt}{}
\titleformat{\subsection}{\normalsize\bfseries}{}{0pt}{}
\titlespacing*{\section}{0pt}{0.35em}{0.12em}
\titlespacing*{\subsection}{0pt}{0.25em}{0.08em}
\pagestyle{fancy}
\fancyhf{}
\fancyhead[C]{\footnotesize Lit Example Reader \quad Assignment ID: U1L9LE \quad Page \thepage\ of \pageref{LastPage}}
\renewcommand{\headrulewidth}{0.5pt}

\newenvironment{playpassage}{\begin{quote}\footnotesize\setlength{\parskip}{0.08em}\setlength{\parindent}{0pt}}{\end{quote}}
\newcommand{\speaker}[1]{\textbf{#1.}}

\begin{document}

\begin{center}
{\LARGE\bfseries Week 9 Lit Example Reader}\par\vspace{0.02em}
{\large\scshape Julius Caesar: When the Reason Assumes the Verdict}\par\vspace{0.06em}
\rule{0.78\textwidth}{0.5pt}
\end{center}

\textbf{Purpose:} Julius Caesar passages for Week 9: test whether a speaker's reason can be accepted without already granting the conclusion.

\textbf{What to watch for:} Restatement, loaded moral labels, and ``must'' language that assumes the verdict still under debate.

\section*{Before the Passages: The Week 9 Logic Tool}

\begin{itemize}[leftmargin=1.3em,itemsep=0.05em,topsep=0.08em,parsep=0.03em]
\item \textbf{Circular reasoning:} the reason depends on already accepting the conclusion.
\item \textbf{Independent support:} evidence that can be granted first, then tested as proof.
\item \textbf{Rewrite:} rebuild the argument so the reason does not smuggle in the verdict.
\end{itemize}

\section*{Passage 1: Act 2, Scene 1 --- Brutus Alone}

\textbf{Context:} Brutus decides Caesar must die. Watch whether his reasons prove future tyranny or assume it.

\begin{playpassage}
\speaker{Brutus} It must be by his death: and for my part,\\
I know no personal cause to spurn at him,\\
But for the general. He would be crown'd:\\
How that might change his nature, there's the question.\\
\ldots\\
The abuse of greatness is when it disjoins\\
Remorse from power: and, to speak truth of Caesar,\\
I have not known when his affections sway'd\\
More than his reason. But 'tis a common proof,\\
That lowliness is young ambition's ladder\ldots\\
So Caesar may.\\
Then, lest he may, prevent. And since the quarrel\\
Will bear no colour for the thing he is,\\
Fashion it thus; that what he is, augmented,\\
Would run to these and these extremities:\\
And therefore think him as a serpent's egg\\
Which, hatch'd, would, as his kind, grow mischievous,\\
And kill him in the shell.
\end{playpassage}

\subsection*{Reader Notes}

Here is your most assisted example. Brutus admits little personal cause and that Caesar's present affections have not outrun reason. The pressure moves to what Caesar \textit{may} become. Week 9 asks whether ``It must be by his death'' rests on independent proof of tyranny, or whether the necessity of killing is already assumed. Notice ``Fashion it thus'': Brutus is shaping the argument toward a verdict.

\textbf{Reading task:} Underline ``It must be by his death.'' Circle present-Caesar vs future-Caesar lines. In the margin, write whether the strongest reason is independent evidence or assumed future danger.

\newpage

\section*{Passage 2: Act 3, Scene 2 --- Brutus to the People}

\textbf{Context:} Brutus explains the killing publicly. Watch whether \textit{ambition} and love of Rome function as independent proof or as labels that assume the verdict.

\begin{playpassage}
\speaker{Brutus} If then that friend demand why Brutus rose against Caesar,\\
this is my answer:---Not that I loved Caesar less, but that I\\
loved Rome more. Had you rather Caesar were living and die all\\
slaves, than that Caesar were dead, to live all free men?\\
As Caesar loved me, I weep for him; as he was fortunate, I rejoice\\
at it; as he was valiant, I honour him: but, as he was ambitious,\\
I slew him.\\
Who is here so base that would be a bondman? If any, speak;\\
for him have I offended. Who is here so rude that would not be a Roman?\\
If any, speak; for him have I offended.
\end{playpassage}

\subsection*{Reader Notes}

This passage shifts more work to you. Brutus pairs ambition with death and pairs dissent with baseness. Mark where a fair hearer would still need proof that Caesar's ambition meant enslavement. Decide whether the either/or question (slaves vs free men) assumes the danger that needed to be shown.

\textbf{Reading task:} Box \textit{ambitious}. Underline the slaves/free men alternative. Note one place the speech may restate or assume the conclusion rather than prove it. Save the full rewrite for the worksheet.

\section*{How to Use These Passages on the Worksheet}

The worksheet asks for circularity diagnosis and repair, not a full history of Rome. Ask: can the reason be accepted before the conclusion?

\end{document}
